\section{Objetivos}

Considerando o plano de trabalho aprovado no âmbito do Programa Institucional de Bolsas de Iniciação Científica (PIBIC/FACEPE), o presente relatório parcial tem como objetivo apresentar o progresso alcançado durante os dois primeiros trimestres de execução do projeto, período dedicado principalmente à consolidação dos fundamentos teóricos, domínio das ferramentas computacionais e implementação inicial do algoritmo Variational Quantum Eigensolver (VQE).

\subsection{Objetivo Geral}

Se tem como objetivo geral, desenvolver e implementar etapas iniciais do algoritmo Variational Quantum Eigensolver (VQE) aplicadas à simulação de sistemas moleculares simples, estabelecendo a infraestrutura computacional e metodológica necessária para a simulação quântica de materiais relevantes para baterias.

\subsection{Objetivos Específicos}

\begin{itemize}

\item Assimilar os fundamentos teóricos de computação quântica, química quântica e algoritmos variacionais aplicados à simulação molecular;

\item Dominar o uso de ferramentas e plataformas computacionais Ket, QuBox e Qiskit para construção e simulação de circuitos quânticos;

\item Implementar o algoritmo VQE para moléculas modelo, com ênfase inicial nos sistemas H$_2$ e LiH;

\item Realizar a validação preliminar dos resultados obtidos por meio de comparação com métodos clássicos de química computacional;

\item Estruturar e documentar o fluxo de desenvolvimento computacional em repositório versionado, garantindo reprodutibilidade científica;

\item Disponibilizar publicamente os códigos e notebooks desenvolvidos durante a execução das atividades iniciais do projeto;

\item Consolidar os resultados obtidos na forma do presente relatório parcial de atividades.

\end{itemize}